\chapter{الكسور الأولى}\label{ch:g4:fractions}

نصف تفاحة، وربع ساعة، وثلاثة أرباع الصف:
الكسور تسمّي أجزاء الأشياء. يقدّم هذا الفصل
الكتابة $\frac{a}{b}$، ويقرأ الكسور على الصور ومستقيم الأعداد،
ويقارن أبسطها. (والكسور تكبر في
\cref{ch:g6:fractions}.)

\section{تسمية الأجزاء}

\begin{definition}[الكسر]\label{def:g4:fractions:def}
اقطع الوحدة إلى أجزاء متساوية وخذ بعضها --- فذلك
\emph{كسر}\index{كسر}:
\[
\frac{3}{4}
\quad
\begin{array}{l}
3 \text{: عدد الأجزاء المأخوذة (وهو \emph{البسط}\index{بسط})؛}\\
4 \text{: عدد الأجزاء المتساوية التي قُطعت إليها الوحدة (وهو
\emph{المقام}\index{مقام}).}
\end{array}
\]
والكسر $\frac12$ هو \emph{النصف}، و $\frac13$ \emph{الثلث}،
و $\frac14$ \emph{الربع}، و $\frac{1}{10}$ \emph{العُشر}.
\end{definition}

\begin{omfigure}
\begin{tikzpicture}[scale=0.78]
  % half circle
  \draw[thick] (0,0) circle (1.0);
  \fill[omDef!40] (0,0) -- (1,0) arc (0:180:1) -- cycle;
  \draw[thick] (-1,0) -- (1,0);
  \node[below, font=\small] at (0,-1.2) {$\frac12$};
  % quarters
  \begin{scope}[xshift=3.6cm]
  \draw[thick] (0,0) circle (1.0);
  \fill[omThm!40] (0,0) -- (1,0) arc (0:90:1) -- cycle;
  \fill[omThm!40] (0,0) -- (0,1) arc (90:180:1) -- cycle;
  \fill[omThm!40] (0,0) -- (-1,0) arc (180:270:1) -- cycle;
  \draw[thick] (-1,0) -- (1,0);
  \draw[thick] (0,-1) -- (0,1);
  \node[below, font=\small] at (0,-1.2) {$\frac34$};
  \end{scope}
  % thirds bar
  \begin{scope}[xshift=6.6cm, yshift=-0.35cm]
  \foreach \i in {0,1,2} \draw[omExo] (\i*1.3,0) rectangle
    (\i*1.3+1.3,0.7);
  \fill[omMeth!40] (0.04,0.04) rectangle (1.26,0.66);
  \draw[thick] (0,0) rectangle (3.9,0.7);
  \node[below, font=\small] at (1.95,-0.15) {$\frac13$};
  \end{scope}
  % tenths bar
  \begin{scope}[xshift=11.6cm, yshift=-0.35cm]
  \foreach \i in {0,...,9} \draw[omExo] (\i*0.42,0) rectangle
    (\i*0.42+0.42,0.7);
  \foreach \i in {0,...,6} \fill[omProp!40] (\i*0.42+0.03,0.04)
    rectangle (\i*0.42+0.39,0.66);
  \draw[thick] (0,0) rectangle (4.2,0.7);
  \node[below, font=\small] at (2.1,-0.15) {$\frac{7}{10}$};
  \end{scope}
\end{tikzpicture}

{\small قراءة الكسور على الصور. ويجب أن تكون الأجزاء
\emph{متساوية}: فثلاثة أجزاء غير متساوية لا تصنع أثلاثًا!}
\end{omfigure}

\begin{example}[كسور مجموعة]\label{ex:g4:fractions:collection}
$\frac13$ من $12$ كرة: اقسم الكرات $12$ إلى $3$ مجموعات
متساوية من $4$؛ فمجموعة واحدة هي $\frac13$، إذن $\frac13$ من $12$ هو $4$
--- و $\frac23$ من $12$ مجموعتان، أي $8$.
\end{example}

\section{الكسور على مستقيم الأعداد}

\begin{method}[وضع $\frac{a}{b}$]\label{met:g4:fractions:line}
اقطع كل وحدة من المستقيم إلى $b$ خطوات متساوية؛ ثم سِر من $0$
مقدار $a$ خطوة. وإن كان $a$ أكبر من $b$، تجاوز السير $1$:
فالكسور يمكن أن تكون أكبر من وحدة كاملة!
\end{method}

\begin{omfigure}
\begin{tikzpicture}[scale=1.6]
  \draw[-Stealth] (-0.2,0) -- (6.9,0);
  \foreach \i in {0,4,8,12,16,20,24}
    \draw ({\i*0.25},-0.09) -- ({\i*0.25},0.09);
  \foreach \i in {0,...,24} \draw ({\i*0.25},-0.05) -- ({\i*0.25},0.05);
  \foreach \x/\l in {0/0, 1/1, 2/2, 3/3, 4/4, 5/5, 6/6}
    \node[below, font=\small] at (\x,-0.1) {$\l$};
  \fill[omDef] (0.5,0) circle (1.6pt)
    node[above, font=\small] {$\frac24$};
  \fill[omThm] (1.25,0) circle (1.6pt)
    node[above, font=\small] {$\frac54$};
  \fill[omMeth] (2.75,0) circle (1.6pt)
    node[above, font=\small] {$\frac{11}{4}$};
\end{tikzpicture}

{\small المستقيم مقطوعًا أرباعًا: $\frac24$ يقع على النقطة
نفسها التي يقع عليها $\frac12$؛ و $\frac54$ يزيد ربعًا على $1$؛
و $\frac{11}{4}$ هو $2$ و $\frac34$.}
\end{omfigure}

\begin{example}[أعداد صحيحة مختبئة في الكسور]\label{ex:g4:fractions:whole}
$\frac44 = 1$ (فأربعة أرباع تصنع الوحدة)؛ و $\frac82 = 4$؛
و $\frac{12}{3} = 4$. \omterm{def:g4:fractions:def}{والكسر} $\frac74$ هو $1 + \frac34$: وحدة كاملة
وثلاثة أرباع فوقها.
\end{example}

\section{مقارنة الكسور البسيطة وجمعها}

\begin{proposition}[المقام نفسه]\label{prop:g4:fractions:compare}
عند تساوي \omterm{def:g4:fractions:def}{المقام}، تكون الأجزاء متساوية الحجم --- فيكفي أن
تعدّها:
\[
\frac{2}{8} < \frac{5}{8},
\qquad
\frac{3}{8} + \frac{4}{8} = \frac{7}{8}.
\]
والمقارنة بالوحدة سهلة أيضًا: $\frac{5}{8} < 1 < \frac{9}{8}$
(أقلّ من $8$ أثمان، ثم أكثر منها).
\end{proposition}

\begin{proof}[بيان بالصورة]
تلوين $3$ أثمان ثم $4$ أثمان أخرى من الشريط نفسه يلوّن
$7$ أثمان في المجموع.
\end{proof}

\begin{example}[مقامان مختلفان، ونقطة واحدة]\label{ex:g4:fractions:equal}
على مستقيم الأعداد، يقع $\frac12$ و $\frac24$ و $\frac{5}{10}$
كلها على النقطة نفسها: فهي ثلاثة أسماء للعدد نفسه. وقطع
كل نصف إلى نصفين يصنع أرباعًا؛ وإلى خمسة، أعشارًا. (والقاعدة
العامة في \cref{ch:g6:fractions}.)
\end{example}

\begin{example}[أرباع الساعة]\label{ex:g4:fractions:hour}
في الساعة $60$ دقيقة، إذن ربع الساعة هو
$60 \div 4 = 15$ دقيقة، وثلاثة أرباع الساعة هي
$45$ دقيقة. و«الثانية والنصف» هي $\frac12$ ساعة بعد الساعة
الثانية.
\end{example}

\section{تمارين}

\begin{exercise}[$\star$]\label{exo:g4:fractions:1}
اُكتب \omterm{def:g4:fractions:def}{الكسر} الظاهر: بيتزا مقطوعة إلى $6$ وبقيت منها $5$ قطع؛
ولوح شوكولاتة من $8$ مربعات أُكل منها $3$ (\omterm{def:g4:fractions:def}{الكسر} المأكول)؛
وعلم مقسوم إلى $3$ أشرطة عمودية متساوية، لوّن منها $1$.
\end{exercise}

\begin{exercise}[$\star$]\label{exo:g4:fractions:2}
اُرسم شريطًا مقطوعًا إلى $5$ أجزاء متساوية ولوّن $\frac35$. ثم لوّن
$\frac{2}{5}$ من شريط آخر مساوٍ له. فأي التلوينين أكبر؟
\end{exercise}

\begin{exercise}[$\star$]\label{exo:g4:fractions:3}
احسب: $\frac12$ من $18$؛ \quad $\frac14$ من $20$؛ \quad
$\frac13$ من $21$؛ \quad $\frac34$ من $20$ (ثلاث مجموعات من
الربع).
\end{exercise}

\begin{exercise}[$\star$]\label{exo:g4:fractions:4}
ضع على مستقيم مقطوع أثلاثًا: $\frac13$؛ \ $\frac33$؛ \
$\frac53$؛ \ $2$. فأي \omterm{def:g4:fractions:def}{كسر} يقع على $2$ بالضبط؟
\end{exercise}

\begin{exercise}[$\star$]\label{exo:g4:fractions:5}
اُنقل وأكمل بالإشارة $<$ أو $>$ أو $=$:
\[
\frac38 \;?\; \frac68, \qquad
\frac55 \;?\; 1, \qquad
\frac74 \;?\; 1, \qquad
\frac12 \;?\; \frac24 .
\]
\end{exercise}

\begin{exercise}[$\star$]\label{exo:g4:fractions:6}
احسب: $\frac26 + \frac36$؛ \quad $\frac58 - \frac28$؛ \quad
$\frac14 + \frac24$؛ \quad $1 - \frac13$ (وكم ثلثًا يصنع
$1$؟).
\end{exercise}

\begin{exercise}[$\star$]\label{exo:g4:fractions:7}
اُكتب على صورة عدد صحيح زائد \omterm{def:g4:fractions:def}{كسر} أصغر من $1$:
$\frac54$؛ \quad $\frac73$؛ \quad $\frac{9}{2}$.
\end{exercise}

\begin{exercise}[$\star$]\label{exo:g4:fractions:8}
كم دقيقة في: نصف ساعة؛ وربع ساعة؛ وثلاثة
أرباع الساعة؛ و $\frac13$ الساعة؟
\end{exercise}

\begin{exercise}[$\star$]\label{exo:g4:fractions:9}
في صفّ من $24$ تلميذًا، $\frac14$ يلبسون نظارات. فكم تلميذًا
يلبس نظارات؟ وكم تلميذًا لا يلبسها؟
\end{exercise}

\begin{exercise}[$\star\star$]\label{exo:g4:fractions:10}
أكلت نور $\frac38$ من بيتزا وأكل سامي $\frac48$ من البيتزا
نفسها. فأي \omterm{def:g4:fractions:def}{كسر} أكلا معًا؟ وأي \omterm{def:g4:fractions:def}{كسر} بقي؟
وأيّهما أكل أكثر؟
\end{exercise}

\begin{exercise}[$\star\star$]\label{exo:g4:fractions:11}
صواب أم خطأ، مع صورة أو مثال: «يمكن أن يكون $\frac12$ من
بيتزا صغيرة أقلّ من $\frac14$ من بيتزا كبيرة جدًّا».
وما الذي يجب أن نعرفه دائمًا قبل مقارنة كسرين من
شيء ما؟
\end{exercise}
